% !TeX root = ../../Book2.tex
% 纲要标题：### D.1 从带对合代数到 \(C^*\)-代数
% 纲要位置：工作记录/计划纲要.md，第 1673 行。

% 纲要要点（供目录核对）：
% * 复代数上的共轭线性对合、Banach 代数及 \(C^*\) 恒等式 \(\|a^*a\|=\|a\|^2\)
% * 例子：\(M_n(\mathbb C)\)、\(C_0(X)\)、\(B(H)\) 与紧算子代数；说明单位元和完备性的不同情况
% * 抽象环、带对合代数和算子代数的附加结构不能省略；Weyl/Ore 微分算子环不因此自动成为 \(C^*\)-代数

\section{从带对合代数到 \texorpdfstring{\(C^*\)}{C*}-代数}
\label{b2:sec:D01}

矩阵既可以相乘，也可以取共轭转置、作用于带内积的空间并计算算子范数。这些附加结构彼此有关，却不能从抽象环的乘法中一并省略。以下讨论复代数；Hilbert 空间的内积沿用\secref{b2:sec:1004}的约定，对第一变量共轭线性、第二变量线性。所有表示均由有界算子给出，除非明确另作说明。

\subsection{对合、Banach 代数与 \texorpdfstring{\(C^*\)}{C*} 恒等式}
\label{b2:subsec:D01:01}

\begin{definition}\label{b2:def:banach-star-cstar}
设 \(A\) 为复结合代数。若映射 \(a\mapsto a^*\) 对任意 \(a,b\in A\)、\(\lambda\in\mathbb C\) 满足
\[
(a+b)^*=a^*+b^*,\quad (\lambda a)^*=\bar\lambda a^*,\quad
(ab)^*=b^*a^*,\quad (a^*)^*=a.
\]
则称它为 \(A\) 上的\term[gonge-xianxing-duihe]{共轭线性对合}[conjugate-linear involution]。若 \(A\) 带有使其完备、并满足 \(\|ab\|\le\|a\|\cdot\|b\|\) 的范数，则称 \(A\) 为\term[Banach-daishu]{Banach 代数}[Banach algebra]。若 \(A\) 还是带有上述对合的 Banach 代数，并对任意 \(a\in A\) 满足
\begin{equation}\label{b2:eq:cstar-identity}
\|a^*a\|=\|a\|^2\qquad(a\in A),
\end{equation}
则称 \(A\) 为 \term[Cxing-daishu]{\(C^*\)-代数}[C*-algebra]。这里不默认存在单位元。若两个带对合的复代数之间的映射保持乘法、复线性结构和对合，则称它为 \term[xing-tongtai]{\(*\)-同态}[*-homomorphism]；需要保持单位元时另行注明。
\end{definition}

由 \(\|a\|^2=\|a^*a\|\le\|a^*\|\cdot\|a\|\)，对 \(a\ne0\) 消去范数得到 \(\|a\|\le\|a^*\|\)；再把 \(a\) 换为 \(a^*\)，可得 \(\|a^*\|=\|a\|\)。因此对合自动连续。非零含幺情形还给出 \(1^*=1\)、\(\|1\|=1\)。这里共轭线性的含义尤其要与附录B中固定基域的线性对合区别：复数标量 \(i\) 被送到 \(-i\)。

\subsection{矩阵、函数与算子代数}
\label{b2:subsec:D01:02}

设 \(H\ne0\) 是复 Hilbert 空间。全部有界线性算子组成 \(B(H)\)，乘法为复合，对合为伴随，范数为
\[
\|T\|=\sup_{\|\xi\|\le1}\|T\xi\|.
\]
\symidx{\(B(H)\)：Hilbert 空间上的有界线性算子代数}
有界算子空间的完备性是这里所需的分析学基础。伴随满足 \(\langle T\xi,\eta\rangle=\langle\xi,T^*\eta\rangle\)，并由内积的\term[Cauchy-Schwarz-budengshi]{Cauchy--Schwarz 不等式}[Cauchy--Schwarz inequality]及单位向量上的上确界得到 \(\|T^*\|=\|T\|\)。

\begin{proposition}\label{b2:prop:operator-cstar}
设 \(H\ne0\) 为复 Hilbert 空间，\(B(H)\) 为其有界复线性算子代数，对合取 Hilbert 空间伴随，则 \(B(H)\) 在算子范数下是含幺 \(C^*\)-代数；它的每个范数闭且对伴随封闭的复子代数也是 \(C^*\)-代数，未必含有 \(I_H\)。
\end{proposition}
\begin{proof}
次乘性来自 \(\|ST\xi\|\le\|S\|\cdot\|T\xi\|\)。对单位向量 \(\xi\)，
\[
\|T\xi\|^2=\langle\xi,T^*T\xi\rangle\le\|T^*T\|,
\]
故 \(\|T\|^2\le\|T^*T\|\)。反向由次乘性及 \(\|T^*\|=\|T\|\) 得到。闭子代数继承完备性、对合和范数恒等式，所以也满足定义。
\end{proof}

取 \(H=\mathbb C^n\)，即得 \(M_n(\mathbb C)\)；其范数是 Euclid 内积诱导的算子范数，不能任意改用另一种矩阵范数。例如把算子范数统一乘二，\(E_{11}\) 就不再满足\eqref{b2:eq:cstar-identity}。

\ProseParagraphBreak
设 \(X\) 为局部紧 Hausdorff 空间。记 \(C_0(X)\) 为连续复函数 \(f\) 的代数，要求对每个 \(\varepsilon>0\)，集合 \(\biggl\{\,x\in X:\Bigl|f(x)\Bigr|\ge\varepsilon\,\biggr\}\) 紧。这是“在无穷远处趋于零”的精确定义。对合为逐点复共轭，范数为上确界范数；一致极限保持上述性质，故它完备，而 \(\|\bar f f\|_\infty=\|f\|_\infty^2\)。当 \(X\) 紧时，它就是含有常值函数一的 \(C(X)\)；\(C_0(\mathbb R)\) 则没有单位元。若它有单位 \(e\)，对每个点选一个在该点非零的紧支集连续函数，\(ef=f\) 迫使 \(e\) 处处为一，与趋零条件矛盾。
\symidx{\(C_0(X)\)：在无穷远处趋于零的连续复函数代数}
\symidx{\(C(X)\)：紧 Hausdorff 空间上的连续复函数代数}

\ProseParagraphBreak
有限秩算子组成 \(B(H)\) 的双边 \(*\)-理想；其范数闭包记为 \(\mathcal K(H)\)，称为\term[jinsuanzi-daishu]{紧算子代数}[algebra of compact operators]。在 Hilbert 空间上，这与把有界集送到相对紧集的算子一致。确实，有限秩算子的单位球像在有限维空间中全有界；范数极限仍使单位球像全有界。反过来，若 \(T\) 的单位球像相对紧，取其有限 \(\varepsilon\)-网，让 \(P\) 为网中向量张成空间的正交投影，便有 \(\|T-PT\|\le\varepsilon\)，而 \(PT\) 有限秩。完备空间中全有界集的闭包紧，故两种定义等价，亦见 Blackadar\cite[\nopp I.8.1.1; I.8.1.5]{Blackadar2006OperatorAlgebras}。

若 \(H\) 无限维，任何有限秩 \(F\) 都有一个单位核向量，故 \(\|I-F\|\ge1\)，所以 \(I\notin\mathcal K(H)\)。它也没有别的单位：若 \(e\) 对全部秩一算子作左单位，逐一作用于其像向量就得到 \(e\xi=\xi\) 对所有 \(\xi\)，从而 \(e=I\)。
\symidx{\(\mathcal K(H)\)：Hilbert 空间上的紧算子代数}

\subsection{抽象环、对合代数与算子代数}
\label{b2:subsec:D01:03}

上三角矩阵代数有天然算子范数，却一般不对共轭转置封闭；有限秩算子对伴随封闭，却在无限维时不完备。一个例子缺少对合封闭性，另一个缺少闭性，不能仅因都由算子组成就称它们为 \(C^*\)-代数。

\begin{proposition}\label{b2:prop:bounded-weyl-impossible}
设 \(A\ne0\) 为含幺复 Banach 代数，则不存在元素 \(x,d\in A\) 满足 \(dx-xd=1_A\)。因此，对任意非零复 Hilbert 空间 \(H\)，复 Weyl 代数 \(A_1(\mathbb C)\) 没有取值于 \(B(H)\) 的保幺代数表示。
\end{proposition}
\begin{proof}
若关系成立，归纳展开交换子得到 \(dx^n-x^nd=nx^{n-1}\)。于是
\[
n\|x^{n-1}\|\le 2\|d\|\cdot\|x^n\|\qquad(n\ge1).
\]
\(d\ne0\)，从 \(\|1\|>0\) 开始迭代，得到
\[
\|x\|^n\ge\|x^n\|\ge\frac{n!}{\Bigl(2\|d\|\Bigr)^n}\|1\|.
\]
但 \((n!)^{1/n}\) 无界，例如最后至少 \(n/2\) 个因子都不小于 \(n/2\)。取 \(n\) 次方根即得矛盾。
\end{proof}

第20章的求导与乘坐标作用并没有违反此命题：它们首先定义在多项式空间上，若置于通常的函数 Hilbert 空间，便涉及无界算子及共同定义域。处理这类作用需要另外的分析，而不能直接把 Ore 扩张的代数表示当作有界表示。
